% !TeX root = ../../Book2.tex
\OptionalSection{色散衰减与频率局部化}{b2:sec:1104}

能量守恒不排斥某些逐点振幅随时间减小：相同的平方能量可以分布在越来越大的区域上。为了看清波动传播中的振荡抵消，先考察三维、固定频带上的衰减估计。频带和空间正则性是结论的一部分；不能删掉频率限制而仍保持同一个 \(L^1\) 数据范数。

\subsection{三维频带上的 \texorpdfstring{\(L^1\to L^\infty\)}{L¹→L∞} 衰减}
\label{b2:subsec:110401}

取 \(c>0\)、\(t\in\mathbb R\)、\(h\in L^2(\mathbb R^3;\mathbb C)\)，定义一支波传播算子
\[
 U_c(t)h\coloneqq\mathcal F_2^{-1}
 \left(e^{2\pi\mathrm{i}ct|\xi|}\widehat h(\xi)\right).
\]
它在 \(L^2\) 上等距；另一支以 \(-t\) 代替 \(t\)。取固定 \(\chi\in C_c^\infty((1/2,2))\)，用径向乘子
\[
 \widehat{P_Nh}(\xi)=\chi(|\xi|/N)\widehat h(\xi),\quad N>0,
\]
把频率局部化到大小为 \(N\) 的环带。这里称为\term[pinlv-jubuhua]{频率局部化}[frequency localization]，其空间作用是卷积，并不是把函数截断到某个空间区域。

\Needspace{6\baselineskip}
\begin{theorem}[三维频带衰减]\label{b2:thm:11-localized-decay}
设 \(c>0\)、\(\chi\in C_c^\infty((1/2,2))\)，\(U_c(t)P_N\) 是符号 \(e^{2\pi\mathrm ict|\xi|}\chi(|\xi|/N)\) 的 Fourier 乘子。存在仅依赖 \(c\)、\(\chi\) 的常数 \(C\)，使任意 \(N>0\)、\(t\in\mathbb R\)、\(h\in L^1(\mathbb R^3;\mathbb C)\) 满足
\begin{equation}\label{b2:eq:11-localized-decay}
 \|U_c(t)P_Nh\|_\infty
 \leq C N^3(1+N|t|)^{-1}\|h\|_1.
\end{equation}
此处由光滑紧频带的卷积核定义算子，因而不要求 \(h\in L^2\)。
\end{theorem}
\begin{proof}
先取 \(N=1\)。核为
\[
 K(t,x)=\int_{\mathbb R^3}e^{2\pi\mathrm{i}(x\cdot\xi+ct|\xi|)}
                       \chi(|\xi|)\,\mathrm{d}\xi.
\]
绝对值直接积分给出统一常数界，足以处理 \(|t|\leq1\)。令 \(r=|x|\)。若 \(r>0\)，对角变量积分，利用
\(\int_{\mathbb S^2}e^{2\pi\mathrm{i}r\rho\theta_1}\,\mathrm{d}\sigma
=2\sin(2\pi r\rho)/(r\rho)\)，得到
\[
 K(t,x)=\dfrac1{\mathrm{i}r}\bigl(I(ct+r)-I(ct-r)\bigr),\quad
 I(s)=\int_0^\infty\rho\chi(\rho)e^{2\pi\mathrm{i}s\rho}\,\mathrm{d}\rho.
\]
当 \(r\geq c|t|/2\) 时，\(I\) 的绝对值有统一界，因此 \(|K|\leq C/|t|\)。这个区域包括最难抵消的传播球面 \(r=c|t|\)。

若 \(r<c|t|/2\)，在角积分之前对径向变量操作。相位的径向导数为 \(ct+x\cdot\theta\)，绝对值至少为 \(c|t|/2\)。因 \(\rho^2\chi(\rho)\) 光滑且紧支撑在正半轴内部，两次分部积分没有边界项，给出
\[
 \left|\int_0^\infty\rho^2\chi(\rho)
 e^{2\pi\mathrm{i}\rho(ct+x\cdot\theta)}\,\mathrm{d}\rho\right|
 \leq C|t|^{-2}
\]
并且常数与方向无关。球面积分后仍有同样估计，也覆盖 \(r=0\)。两个区域合起来得到 \(\|K(t)\|_\infty\leq C(1+|t|)^{-1}\)。

一般 \(N\) 时换元 \(\xi=N\eta\)，有 \(K_N(t,x)=N^3K(Nt,Nx)\)。核的上确界遂满足\eqref{b2:eq:11-localized-decay}中的系数。最后对 \(K_N*h\) 使用绝对积分界，得到算子估计。
\end{proof}

对于所有 \(N>0\)、\(t\ne0\)，都有 \(N^3(1+N|t|)^{-1}\leq |t|^{-1}N^2\)。两个频率幂说明，全频率估计必须支付相应的空间正则性代价。具体地，设 \((h_j)_{j\in\mathbb Z}\subset L^1(\mathbb R^3;\mathbb C)\)，其 Fourier 变换满足
\[
 \operatorname{supp}\widehat h_j\subseteq
 \{\xi:2^{j-1}\leq|\xi|\leq2^{j+1}\},
\]
并且
\[
 \sum_j2^{2j}\|h_j\|_1<\infty,
\]
每个 \(\widehat h_j\) 有界且紧支撑，因而也属于 \(L^2\)；由 Fourier 变换的 \(L^1\) 与 \(L^2\) 版本相容，\(h_j\) 也属于 \(L^2\)，所以此前定义的 \(U_c(t)h_j\) 有意义。取 \(\widetilde\chi\in C_c^\infty((1/4,4))\)，使其在 \([1/2,2]\) 上恒等于一。前面核估计的证明只用截断光滑且紧支撑在正半轴内，对这一稍宽环带同样适用；将截断取为 \(\widetilde\chi(|\xi|/2^j)\)，它在 \(\widehat h_j\) 的支集上为一。逐项应用核估计，得到
\[
 \left\|\sum_jU_c(t)h_j\right\|_\infty
 \leq\dfrac C{|t|}\sum_j2^{2j}\|h_j\|_1,\quad t\ne0.
\]
因此对每个固定的 \(t\ne0\)，传播后的级数在 \(L^\infty\) 中绝对收敛。级数的和因而给出由这些频带数据确定的有界函数。

\ProseParagraphBreak

若另有 \(h\in L^2(\mathbb R^3;\mathbb C)\)，且对称部分和 \(H_M=\sum_{|j|\leq M}h_j\) 在 \(L^2\) 中趋于 \(h\)，则上述级数的和确实等于此前的 \(U_c(t)h\)。证明要同时识别两种范数极限。记 \(V_M=\sum_{|j|\leq M}U_c(t)h_j\)，\(v=\lim_M V_M\) 为 \(L^\infty\) 极限。由 \(U_c(t)\) 的线性与 \(L^2\) 等距性，
\[
 V_M=U_c(t)H_M\longrightarrow U_c(t)h\quad\text{于 }L^2.
\]
对任意半径 \(R>0\)，两种收敛给出
\[
 \|v-U_c(t)h\|_{L^2(B_R)}
 \leq |B_R|^{1/2}\|v-V_M\|_\infty
       +\|V_M-U_c(t)h\|_2\longrightarrow0.
\]
故在每个球上两个极限几乎处处相同；对正整数半径的球合并零测例外，便有 \(v=U_c(t)h\) 几乎处处成立于全空间。这里初始部分和的 \(L^2\) 收敛同时控制了高频与低频的截断极限。仅有初始频带和在分布意义下收敛还不足以作这一步；当 \(t\ne0\) 时，符号 \(e^{2\pi\mathrm{i}ct|\xi|}\) 在原点不光滑，不能无条件与任意温和分布相乘。

\subsection{色散}
\label{b2:subsec:110402}

在振荡积分中，不同方向的相位可以互相抵消；传播球面附近保留最大的贡献，而远离它的区域允许更多次分部积分。这种由相位和传播方向使振幅分散的现象，在本节称为\term[sesan]{色散}[dispersion]。波方程的关系式 \(\rho=2\pi c|\xi|\) 对径向频率是线性的，沿同一方向的不同频率以相同速度行进；因此它的衰减机制包含角向展开，不能直接照搬二次相位的热核或 Schrödinger 核计算。

\ProseParagraphBreak

频率代价不是上述证明的偶然损失。取非负非零 \(\chi\)，令 \(h=\mathcal F^{-1}(\chi(|\xi|))\)，再令 \(h_N(x)=N^3h(Nx)\)。这些都是 Schwartz 函数，且 \(\|h_N\|_1=\|h\|_1\)。在固定 \(t>0\) 和 \(|x|=ct\) 处，刚才的角积分公式给出
\[
 U_c(t)h_N(x)=\dfrac{N^2}{\mathrm{i}ct}
                    \bigl(I(2Nct)-I(0)\bigr).
\]
分部积分说明 \(I(2Nct)\to0\)，而 \(I(0)>0\)，故其模以 \(N^2\) 的量级增长。于是不存在对全部 Schwartz 数据成立的固定时间无导数 \(L^1\to L^\infty\) 界；更不能宣称完整波群具有 \(C|t|^{-1}\|h\|_1\) 的三维界。

\ProseParagraphBreak

另一方面，\(U_c(t)\) 在 \(L^2\) 中没有范数衰减。这与频带上的逐点衰减并不冲突：上确界下降时，波可以分布在更大的球壳上。三维壳层的面积按 \(t^2\) 增长，幅度按 \(t^{-1}\) 减小，恰与平方能量保持同一尺度相容。这只是对估计的几何解释；严格结论仍是上面带频带条件的核界与 Plancherel 等式。

\subsection{Strichartz 估计的动机}
\label{b2:subsec:110403}

非线性波方程中的外力依赖解本身，仅有每个时刻的能量界往往不够控制 Duhamel 积分。对区间 \(I\subset\mathbb R\)、可测 \(u:I\times\mathbb R^n\to\mathbb C\) 及 \(1\le q,r\le\infty\)，需要同时记录时间和空间的可积性，例如
\[
 \|u\|_{L^q_tL^r_x(I\times\mathbb R^n)}
 =\left(\int_I\left(\int_{\mathbb R^n}|u(t,x)|^r\,\mathrm{d}x\right)^{q/r}
 \,\mathrm{d}t\right)^{1/q},
\]
有限指数时按此定义，指数为无穷时改用本性上确界。能量给出某些 \(q=\infty\) 的控制，色散则提供随时间变化的空间映射界；连接两者的时空估计属于\term[strichartz-guji]{Strichartz 估计}[Strichartz estimates]的研究范围。

\ProseParagraphBreak

从单个频带的核估计走向时空估计，还需要处理时间积分的端点、频带求和以及数据的正则性。陶哲轩的《Nonlinear Dispersive Equations: Local and Global Analysis》\cite[Section 2.3]{Tao2006Nonlinear}将此类时空估计用于非线性色散方程的局部与整体分析；本卷\appref{b2:app:E}将进一步讨论相应的时空估计。

